%% methodology.tex
%%

%% ==============================
\section{Methodology}
\label{sec:Methodology}
%% ==============================

This section describes the experimental methodology used to evaluate whether cross-head fusion reduces the systematic underprediction reported in the baseline paper. Section~\ref{subsec:synth} specifies the synthetic telemetry generator; Section~\ref{subsec:window} details the forecasting window design; Section~\ref{subsec:transient} describes the transient burst injection mechanism; Section~\ref{subsec:training} presents the training procedure; and Section~\ref{subsec:metrics} defines the evaluation metrics.

\subsection{Synthetic Telemetry Generator}
\label{subsec:synth}

Real multi-tenant cluster telemetry was not available for this project. While the Google Borg traces~\cite{Tirmazi20,Liu12} and IBM Cloud dataset~\cite{CloudAnomalyIBM24} provide large-scale production workload data, they do not include ground-truth labels for future violations or controlled cross-metric dependencies required to isolate the specific failure mode this project targets. Therefore, a synthetic telemetry generator (\texttt{src/data/telemetry\_simulator.py}) was developed to produce controlled, reproducible multivariate time series with known cross-metric burst precursors.

The generator produces sequences of 15 timesteps for four metrics: CPU utilization, memory utilization, disk I/O throughput, and network throughput. Each metric is represented as a normalized value in the range [0, 1]. The 15-step sequence is partitioned into a 10-step \emph{history window} (timesteps 0--9, the only input the models see) and a 5-step \emph{future window} (timesteps 10--14, used only to compute ground-truth labels, never shown to the models). This design makes the task a genuine forecasting problem: the model must predict future violations (timesteps 10--14) based solely on observed history (timesteps 0--9).

\textbf{Design iteration}: An earlier version of this experiment computed labels from the same window the models could observe, which trivially favoured the reactive threshold baseline (since it could ``predict'' violations that had already occurred in the visible window). This was identified as a design flaw during the project and corrected: the current version strictly separates history (input) from future (labels), ensuring that all models -- baseline, per-head, and fused -- are evaluated on a genuine forecasting task with the same information asymmetry.

\subsection{Forecasting Window Design}
\label{subsec:window}

Each of the four metrics has independent None/L1/L2 violation thresholds:
\begin{itemize}
  \item \textbf{None}: $x < 0.6$ (normal operating range)
  \item \textbf{Level 1 (L1)}: $0.6 \leq x < 0.8$ (elevated usage, approaching capacity)
  \item \textbf{Level 2 (L2)}: $x \geq 0.8$ (severe usage, imminent capacity exhaustion)
\end{itemize}

These thresholds are fixed and uniform across all metrics for simplicity. Real cluster environments typically use per-metric, per-node thresholds tuned to hardware characteristics and workload patterns, but uniform thresholds suffice for a controlled experiment focused on cross-metric dependencies rather than threshold calibration.

For each timestep in the future window (10--14), the generator computes a ground-truth violation label (None/L1/L2) per metric by applying these thresholds to the peak value over that timestep's 1-step window. The label used for training and evaluation is the \emph{maximum severity} observed across the 5-step future window for each metric. This design choice reflects the operational reality that predicting \emph{any} future violation in the next 5 steps is valuable, even if the exact timestep is uncertain.

\subsection{Transient Burst Injection}
\label{subsec:transient}

Non-transient samples (60\% of the dataset) are generated as independent Brownian noise per metric, with no cross-metric correlation. These represent normal cluster operation, where metrics fluctuate independently within their normal ranges.

Transient samples (40\% of the dataset) are engineered to contain cross-metric burst precursors:
\begin{enumerate}
  \item A \textbf{primary metric} is randomly selected (CPU, memory, disk, or network).
  \item A \textbf{fast, non-linear ramp} is injected into the primary metric, starting at timestep $t_{\text{primary}} \in [7, 9]$ (late in the history window or early in the future window) and ramping to a peak value $v_{\text{peak}} \in [0.7, 0.95]$ over 2--3 timesteps. This ramp lands mostly or fully in the future window (timesteps 10--14), ensuring the primary metric's violation is a genuine forecasting target.
  \item One or two \textbf{rider metrics} (randomly selected from the remaining three) are injected with a correlated ramp starting 3--6 timesteps \emph{earlier} than the primary metric ($t_{\text{rider}} = t_{\text{primary}} - \delta$, where $\delta \in [3, 6]$). The rider ramps to a lower peak ($v_{\text{rider}} \in [0.5, 0.7]$) and is visible in the history window.
  \item A small amount of Gaussian noise ($\sigma = 0.05$) is added to all metrics to simulate measurement variability.
\end{enumerate}

This construction ensures that predicting the primary metric's future violation (timesteps 10--14) sometimes requires reading a \emph{different} metric's early signal (timesteps 4--9) -- exactly the scenario a strict per-head architecture (where heads never exchange information) struggles with. The rider metric's precursor is intentionally kept below the L1 threshold (peak $< 0.6$) so the rider itself does not trigger a violation, but its temporal correlation with the primary metric is the key predictive signal.

Figure~\ref{fig:transient_example} illustrates a transient sample where a network spike (rider, timesteps 5--8) precedes a CPU spike (primary, timesteps 11--13). The CPU head in a strict per-head architecture cannot leverage the network head's early signal, leading to underprediction of the CPU violation.

\begin{figure}[htp]
\centering
\fbox{\parbox{0.9\textwidth}{\centering \textit{[Diagram: 4-panel time series plot showing CPU, memory, disk, network over 15 timesteps (0--14), with vertical line at timestep 10 separating history (left) from future (right). Network shows rider spike in history window (timesteps 5--8, peak 0.65); CPU shows primary spike in future window (timesteps 11--13, peak 0.85, L2 violation). Memory and disk remain in normal range throughout.]}}}
\caption{Example transient sample with cross-metric burst precursor: network spike (rider, history window) precedes CPU spike (primary, future window).}
\label{fig:transient_example}
\end{figure}

\subsection{Training Procedure}
\label{subsec:training}

Both attention models (MHSA-PerHead and MHSA-Fused) are trained with the Adam optimizer (learning rate 0.001, default $\beta_1 = 0.9$, $\beta_2 = 0.999$) for 30 epochs on 20,000 samples per seed. The dataset is split 80/20 into train and test sets, stratified on the transient flag to ensure both sets contain the same 40\% transient ratio.

The loss function is a summed per-metric cross-entropy loss with inverse-frequency class weights. Let $y_i^{(m)} \in \{0, 1, 2\}$ be the true label for metric $m$ on sample $i$, and let $\hat{p}_i^{(m)}$ be the predicted probability distribution over classes. The loss for metric $m$ is:
\begin{equation}
\mathcal{L}_m = -\sum_{i=1}^{N} w_{y_i^{(m)}} \log \hat{p}_i^{(m)}[y_i^{(m)}]
\end{equation}
where $w_c$ is the inverse-frequency weight for class $c \in \{0, 1, 2\}$ (None, L1, L2):
\begin{equation}
w_c = \frac{\sum_{c'=0}^{2} n_{c'}}{3 \cdot n_c}
\end{equation}
and $n_c$ is the number of training samples with true label $c$ for metric $m$. This weighting counteracts the dominance of the ``none'' class (which typically accounts for 70--80\% of samples) and ensures rare L1/L2 violations are not drowned out during training.

The total loss is the sum over all four metrics:
\begin{equation}
\mathcal{L}_{\text{total}} = \sum_{m=1}^{4} \mathcal{L}_m
\end{equation}

Five random seeds (42, 43, 44, 45, 46) are used to regenerate the dataset and retrain both models. Critically, \textbf{both} the data generator's random number generator (NumPy) \textbf{and} PyTorch's random number generator (used for model weight initialization and DataLoader shuffling) are seeded independently per run. An earlier, unseeded version of this experiment showed a much larger apparent improvement from fusion that did not replicate once training randomness was properly controlled (see Section~\ref{sec:discussion}). The properly-seeded, 5-seed results reported in Section~\ref{sec:Evaluation} are the ones that should be trusted.

\subsection{Evaluation Metrics}
\label{subsec:metrics}

Four metrics are computed over the held-out test set to evaluate model performance:

\begin{itemize}
  \item \textbf{Accuracy}: Per-metric classification accuracy (fraction of correct None/L1/L2 predictions), averaged across the four metrics. This measures overall prediction correctness but does not distinguish between different types of errors.

  \item \textbf{Macro-F1}: Per-metric macro-averaged F1 score (harmonic mean of precision and recall, averaged across the three classes with equal weight), averaged across the four metrics. This balances precision and recall and is more robust to class imbalance than accuracy alone.

  \item \textbf{Transient Violation Recall}: Of the test samples that are both transient \emph{and} have a true future violation (L1 or L2) on a given metric, the fraction where the model predicts \emph{any} violation (L1 or L2, not necessarily the correct severity). This metric isolates the model's ability to detect that a violation is coming during high-load transients, which is the failure mode the baseline paper reports. A model with high transient recall may still underpredict severity (e.g., predicting L1 when the true label is L2), but it at least flags the violation.

  \item \textbf{Transient Underprediction Bias}: mean$(y_{\text{true}} - y_{\text{pred}})$ over the same subset (transient samples with true violations). Positive values indicate the model systematically underpredicts severity (e.g., predicting L1 when the true label is L2, or None when the true label is L1). This is the metric that directly operationalizes the baseline paper's own stated gap: ``systematic underprediction during high-load transients.'' A bias of 0 means no systematic under- or over-prediction; values close to 1 indicate the model is consistently one severity level too low.
\end{itemize}

All metrics are computed per-metric and then averaged across the four metrics (CPU, memory, disk, network) to produce a single scalar per model. The mean and standard deviation across the five seeds are reported to quantify result variability.

All code, datasets, results, and trained model weights are available in the \texttt{mhsa-tdl-framework/} directory for reproducibility.

